\documentclass[11pt,a4paper]{article}
% Common preamble for the hint sheets (\input from each file)
\usepackage[utf8]{inputenc}
\usepackage[T1]{fontenc}
\usepackage{lmodern}
\usepackage{amsmath,amssymb,amsthm,mathtools}
\usepackage[a4paper,margin=2.3cm]{geometry}
\usepackage{enumitem}
\usepackage{xcolor}
\usepackage[hidelinks]{hyperref}
\renewcommand{\leq}{\leqslant}
\renewcommand{\le}{\leqslant}
\renewcommand{\geq}{\geqslant}
\renewcommand{\ge}{\geqslant}
\setlength{\parindent}{0pt}
\setlength{\parskip}{0.35em}
% \hintstitle{Sheet title}{Class work / Homework}
\newcommand{\hintstitle}[2]{\begin{center}{\LARGE\bfseries Hints: #1\par}\vspace{0.5em}{\large #2\par}\end{center}
\noindent\textit{The hints to each problem are cumulative. Read one, try again, and open the next only if you are still stuck.}\par\medskip\hrule\medskip}
\newcommand{\problem}[1]{\par\bigskip{\large\bfseries Problem #1}\par\nopagebreak\smallskip}
\newcommand{\subpart}[1]{\par\smallskip\textbf{(#1)}\par\nopagebreak}
\newcommand{\hint}[1]{\par\smallskip\noindent\textcolor{black!70}{\textbf{Hint #1.}}\enspace}
\newcommand{\note}[1]{\par\smallskip\noindent\textit{Note.} #1\par}
\providecommand{\lt}{<}
\providecommand{\gt}{>}
\providecommand{\tightlist}{\setlength{\itemsep}{0pt}\setlength{\parskip}{0pt}}
\emergencystretch=3em

\begin{document}
\hintstitle{Functional sequences, series and spaces -- 4. Spaces of functions}{Class work}

\problem{1}
\hint{1}
Completeness: a Cauchy sequence in the uniform norm converges uniformly (Cauchy criterion), and the limit is continuous. Separability: polynomials with rational coefficients (Weierstrass), or piecewise linear functions with rational nodes and values (uniform continuity), form a countable dense set.
\hint{2}
Non-compactness of the ball: tents \(\varphi_n\) of height \(1\) supported on disjoint intervals \(\bigl[\frac1{n+1},\frac1n\bigr]\) satisfy \(\lVert\varphi_n-\varphi_m\rVert=1\) for \(n\ne m\), so no subsequence is Cauchy.
\hint{3}
Riesz's lemma: for a closed proper subspace \(Y\) and \(\theta\lt1\), take \(x_0\notin Y\), \(d=\operatorname{dist}(x_0,Y)\gt0\), and \(y_0\in Y\) with \(\lVert x_0-y_0\rVert\le d/\theta\); then \(x=\frac{x_0-y_0}{\lVert x_0-y_0\rVert}\) has \(\operatorname{dist}(x,Y)\ge\theta\). In infinite dimensions apply it repeatedly to the (closed, because finite-dimensional) spans of \(x_1,\dots,x_{n-1}\), getting unit vectors with pairwise distances at least \(\frac12\).

\problem{2}
\hint{1}
Banach: \(d(x_{n+1},x_n)\le q^nd(x_1,x_0)\), so \((x_n)\) is Cauchy, with geometric bound \(d(x_n,x)\le\frac{q^n}{1-q}d(x_1,x_0)\); the limit is fixed by continuity, and two fixed points satisfy \(d(x,y)\le q\,d(x,y)\).
\hint{2}
Picard: on a box \(\lvert x-x_0\rvert\le a\), \(\lvert y-y_0\rvert\le b\) let \(\lvert f\rvert\le M\) and \(L\) be the Lipschitz constant. On the closed subset \(\{y\in C[x_0-h,x_0+h]:\lvert y-y_0\rvert\le b\}\) the map \(Ty(x)=y_0+\int_{x_0}^xf(t,y(t))\,dt\) maps the set into itself if \(hM\le b\), and is a contraction if \(hL\lt1\). Fixed points of \(T\) are exactly the solutions (fundamental theorem of calculus).
\hint{3}
Without Lipschitz: \(y'=3y^{2/3}\), \(y(0)=0\) has the solutions \(y\equiv0\) and \(y=x^3\) (and many more glued ones).

\problem{3}
\hint{1}
Use the Faber--Schauder system: \(\varphi_0=1\), \(\varphi_1=x\), and then, for \(k\ge0\) and \(0\le j\lt2^k\), the tent of height \(1\) supported on \(\bigl[\frac j{2^k},\frac{j+1}{2^k}\bigr]\) with peak at the midpoint, ordered level by level.
\hint{2}
Choose the coefficients so that after finishing level \(k\) the partial sum is the piecewise linear interpolant \(L_kf\) of \(f\) at the dyadic points \(\frac i{2^{k+1}}\): the coefficient of a new tent is \(f\)(midpoint) minus the average of \(f\) at the two endpoints. Uniform continuity gives \(\lVert f-L_kf\rVert\le2\,\omega(f,2^{-k-1})\to0\).
\hint{3}
Partial sums inside a level are again piecewise linear interpolants of \(f\) at a set of nodes of mesh at most \(2^{-k}\), so the same estimate holds for every partial sum, not just the block sums. Uniqueness of the coefficients: evaluating a uniformly convergent expansion at the dyadic points recovers them one by one.

\problem{4}
\hint{1}
Non-completeness: continuous \(f_n\) equal to \(0\) on \([0,\frac12]\), to \(1\) on \([\frac12+\frac1n,1]\) and linear in between are \(\rho\)-Cauchy. If \(\rho(f_n,g)\to0\) with \(g\) continuous, then \(\int_0^1\lvert g-\mathbf 1_{(1/2,1]}\rvert=0\), which forces \(g=0\) on \([0,\frac12)\) and \(g=1\) on \((\frac12,1]\), impossible.
\hint{2}
Build the completion: \(W\) is the set of \(\rho\)-Cauchy sequences modulo \((f_n)\sim(g_n)\iff\rho(f_n,g_n)\to0\), with \(\tilde\rho\bigl([f_n],[g_n]\bigr)=\lim\rho(f_n,g_n)\) (the limit exists because \(\lvert\rho(f_n,g_n)-\rho(f_m,g_m)\rvert\le\rho(f_n,f_m)+\rho(g_n,g_m)\)). Addition and scalar multiplication are termwise and well defined because \(\rho\) comes from a norm.
\hint{3}
The embedding is \(f\mapsto[(f,f,\dots)]\); it is an isometry, and its image is dense because \([(f_n)]\) is the \(\tilde\rho\)-limit of the constant sequences \([(f_k,f_k,\dots)]\) as \(k\to\infty\). For completeness of \(W\), approximate a \(\tilde\rho\)-Cauchy sequence \(w_k\) by elements \(f_k\) of the image with \(\tilde\rho(w_k,f_k)\lt\frac1k\); then \((f_k)\) is \(\rho\)-Cauchy and its class is the limit. (Concretely \(W\) is \(L^1[0,1]\).)

\problem{5}
\hint{1}
Let \(c=T(1)\). Since \(T\) commutes with differentiation, \(c'=T(1')=T(0)=0\), so \(c\) is a constant function.
\hint{2}
Fix a point \(a\). By Hadamard's lemma every \(f\in C^\infty(\mathbb R)\) can be written \(f(x)=f(a)+(x-a)g(x)\) with \(g(x)=\int_0^1f'(a+t(x-a))\,dt\) smooth. Apply \(T\), using linearity and the commutation with multiplication by \(x\) (hence by \(x-a\)): \(Tf=cf(a)+(x-a)\,Tg\). Evaluating at \(x=a\) gives \((Tf)(a)=cf(a)\).
\end{document}
